\documentclass[10pt,a4paper]{amsart}
\usepackage[margin=19mm]{geometry}
\usepackage{amsmath,amssymb,mathtools}
\usepackage[scr=rsfs,cal=euler]{mathalfa}
\usepackage{xfrac}
\usepackage[unicode,hidelinks,bookmarks=false]{hyperref}
\newcommand{\ie}{i.e.\ }
\newcommand{\cf}{cf.\ }
\newcommand{\resp}{respectively\ }
\newcommand{\Subsection}{\subsection}
\providecommand{\nobreakdash}{\nobreak\hbox{-}}
\setlength{\emergencystretch}{2em}
\multlinegap0pt
\usepackage{float}
\usepackage{amssymb}
\usepackage{mathtools}
\usepackage[shortlabels]{enumitem}
\newtheorem{theorem}{Theorem}
\title{\textbf{Combinatorial metaplexes and centrality indices for identifying higher-order interactions}}
\author{Hiren J. Dhameliya, Udit Raj, Sudeepto Bhattacharya}
\date{Source: arXiv:2602.23383v2 (CC BY 4.0)}
\begin{document}
\maketitle

\noindent\emph{Source and licence.} \textbf{Combinatorial metaplexes and centrality indices for identifying higher-order interactions}. Version arXiv:2602.23383v2 (CC BY 4.0), distributed by arXiv under the Creative Commons Attribution 4.0 International licence, http://creativecommons.org/licenses/by/4.0/, as recorded in the arXiv licence metadata for this exact version and shown on the abstract page https://arxiv.org/abs/2602.23383v2. Full licence notice: \texttt{licenses/arxiv-2602.23383-CC-BY-4.0.txt}.\par\smallskip

\noindent\textit{Editorial scope.} Counted unit: the article's theorem theorem (source0.tex lines 853--894). The article's complete original proof of it is delivered with it (source0.tex lines 896--1016, one \texttt{proof} environment of 121 lines). No mathematics is written, repaired or re-derived: the statement and the proof are the article's own lines, in the article's own notation, and no retained line is substituted or re-flowed.  No further source-local context is needed: every cross-reference of every retained line resolves inside the two delivered spans.  Nothing was omitted from any retained span, so the package carries no omission record.  No retained line contains a citation, so the wrapper carries no bibliography and no bibliography entry.  No retained line contains a figure environment, an image-inclusion command or drawing code, so no image is delivered and the package declares no asset dependency.  Editorial formatting only, in addition: an AMS \texttt{amsart} preamble that restates the article's own declarations and macros used by the retained lines, namely the environments \texttt{theorem} on separate counters, together with the standard packages those lines need.  The retained environments print in this document as Theorem 1 in order of appearance, whereas the article prints the same items with the numbering of its own full document.  The complete native source of the article as distributed in the arXiv e-print for this exact version is bundled unchanged as \texttt{sources/195395/source0.tex}.
\begin{theorem}
Let $1 \le q \le \dim(\Delta)-1$. 
Let $a_q$ and $a_{q+1}$ be the fractional weight contribution maps. 
Then their composition map is
\[
a_{q+1} \circ a_q : S_{q-1} \times S_{q+1} 
\longrightarrow \mathbb{Q} \setminus \mathbb{Q}^-,
\]
\[
(\sigma^{(q-1)}, \sigma^{(q+1)}) 
\longmapsto 
(a_{q+1} \circ a_q)(\sigma^{(q-1)}, \sigma^{(q+1)}),
\]
where
\[
(a_{q+1} \circ a_q)(\sigma^{(q-1)}, \sigma^{(q+1)})
=
\sum_{\sigma^{(q)} \in S_q}
a_{q+1}(\sigma^{(q)}, \sigma^{(q+1)})
a_q(\sigma^{(q-1)}, \sigma^{(q)}).
\]

The map $a_{q+1} \circ a_q$ satisfies the following conditions:

\begin{itemize}
    \item[(i)] If $\sigma^{(q-1)} \not\subset \sigma^{(q+1)}$, then
    \(
    (a_{q+1} \circ a_q)(\sigma^{(q-1)}, \sigma^{(q+1)}) = 0.
    \)

    \item[(ii)] If $\sigma^{(q-1)} \subset \sigma^{(q+1)}$, then
    \(
    0 < (a_{q+1} \circ a_q)(\sigma^{(q-1)}, \sigma^{(q+1)}) \le 1.
    \)

    \item[(iii)] For each $\sigma^{(q-1)} \in S_{q-1}$ if $\exists \sigma^{(q+1)} \in S_{q+1}$ such that $\sigma^{(q-1)} \subset \sigma^{(q+1)}$, then
    \[
    \sum_{\sigma^{(q+1)} \in S_{q+1}}
    (a_{q+1} \circ a_q)(\sigma^{(q-1)}, \sigma^{(q+1)}) = 1.
    \]
\end{itemize}
\end{theorem}

\begin{proof}

Since
\[
a_q(\sigma^{(q-1)}, \sigma^{(q)}) 
\in \mathbb{Q} \setminus \mathbb{Q}^-,
\qquad
a_{q+1}(\sigma^{(q)}, \sigma^{(q+1)}) 
\in \mathbb{Q} \setminus \mathbb{Q}^-,
\]
their product belongs to $\mathbb{Q} \setminus \mathbb{Q}^-$, and hence
\[
(a_{q+1} \circ a_q)(\sigma^{(q-1)}, \sigma^{(q+1)})
\in \mathbb{Q} \setminus \mathbb{Q}^-.
\]

\medskip

\noindent
\text{Proof of (i):}
If $\sigma^{(q-1)} \not\subset \sigma^{(q+1)}$, then there does not exist 
$\sigma^{(q)} \in S_q$ such that
\[
\sigma^{(q-1)} \subset \sigma^{(q)} 
\subset \sigma^{(q+1)}.
\]
Hence for every $\sigma^{(q)} \in S_q$, either
\[
a_q(\sigma^{(q-1)}, \sigma^{(q)}) = 0
\quad \text{or} \quad
a_{q+1}(\sigma^{(q)}, \sigma^{(q+1)}) = 0.
\]
Therefore every summand in
\[
(a_{q+1} \circ a_q)(\sigma^{(q-1)}, \sigma^{(q+1)})
=
\sum_{\sigma^{(q)} \in S_q}
a_{q+1}(\sigma^{(q)}, \sigma^{(q+1)})
a_q(\sigma^{(q-1)}, \sigma^{(q)})
\]
is zero, and thus
\[
(a_{q+1} \circ a_q)(\sigma^{(q-1)}, \sigma^{(q+1)}) = 0.
\]

\medskip

\noindent
\text{Proof of (ii):}
If $\sigma^{(q-1)} \subset \sigma^{(q+1)}$, then there exists at least one 
$\sigma^{(q)} \in S_q$ such that
\[
\sigma^{(q-1)} \subset \sigma^{(q)} 
\subset \sigma^{(q+1)}.
\]
For such $\sigma^{(q)}$,
\[
a_q(\sigma^{(q-1)}, \sigma^{(q)}) > 0
\quad \text{and} \quad
a_{q+1}(\sigma^{(q)}, \sigma^{(q+1)}) > 0.
\]
Hence at least one summand in the defining sum is strictly positive,
and therefore
\[
0 < (a_{q+1} \circ a_q)(\sigma^{(q-1)}, \sigma^{(q+1)}).
\]
Since 
\[
0 \le a_{q+1}(\sigma^{(q)}, \sigma^{(q+1)}) \le 1
\quad \text{and} \quad
a_q(\sigma^{(q-1)}, \sigma^{(q)}) \ge 0,
\]
we have
\[
a_{q+1}(\sigma^{(q)}, \sigma^{(q+1)})
\, a_q(\sigma^{(q-1)}, \sigma^{(q)})
\le
a_q(\sigma^{(q-1)}, \sigma^{(q)}).
\]
Summing over $\sigma^{(q)} \in S_q$ gives
\begin{align*}
(a_{q+1} \circ a_q)(\sigma^{(q-1)}, \sigma^{(q+1)})
&=
\sum_{\sigma^{(q)} \in S_q}
a_{q+1}(\sigma^{(q)}, \sigma^{(q+1)})
a_q(\sigma^{(q-1)}, \sigma^{(q)}) \\
&\le
\sum_{\sigma^{(q)} \in S_q}
a_q(\sigma^{(q-1)}, \sigma^{(q)}) \\
&= 1.
\end{align*}

\medskip

\noindent
\text{Proof of (iii):}
Let $\sigma^{(q-1)} \in S_{q-1}$ such that there exists
$\sigma^{(q+1)} \in S_{q+1}$ with
$\sigma^{(q-1)} \subset \sigma^{(q+1)}$.
Then
\begin{align*}
\sum_{\sigma^{(q+1)} \in S_{q+1}}
(a_{q+1} \circ a_q)(\sigma^{(q-1)}, \sigma^{(q+1)})
&=
\sum_{\sigma^{(q+1)} \in S_{q+1}}
\sum_{\sigma^{(q)} \in S_q}
a_{q+1}(\sigma^{(q)}, \sigma^{(q+1)})
a_q(\sigma^{(q-1)}, \sigma^{(q)}) \notag \\
&=
\sum_{\sigma^{(q)} \in S_q}
\left(
\sum_{\sigma^{(q+1)} \in S_{q+1}}
a_{q+1}(\sigma^{(q)}, \sigma^{(q+1)})
\right)
a_q(\sigma^{(q-1)}, \sigma^{(q)}) \notag \\
&=
\sum_{\sigma^{(q)} \in S_q}
a_q(\sigma^{(q-1)}, \sigma^{(q)}) \notag \\
&= 1.
\end{align*}
\end{proof}

\end{document}
